% ==============================================================================
% Phân mục 05 [TRỌNG TÂM NHIỆM VỤ 3.2]: Hệ thống Đa tác tử & Chống ảo giác
% ==============================================================================
\section{Hiện Thực Hệ Thống Đa Tác Tử và Cơ Chế Neo Bằng Chứng Chống Ảo Giác (Multi-Agent RAG \& Evidence Grounding)}
\label{sec:multi_agent_rag}

Một trong những yêu cầu khắt khe nhất của y học chứng cứ (\textit{Evidence-Based Medicine}) là mọi gợi ý chẩn đoán phân biệt phải truy nguyên được về hai nguồn căn cứ độc lập: (1) triệu chứng cụ thể ghi nhận trên chính bệnh nhân đó, và (2) điều khoản tương ứng trong hướng dẫn thực hành lâm sàng chính thống. Để đáp ứng yêu cầu này mà không rơi vào bẫy ``ảo giác ngôn ngữ'' (\textit{LLM Hallucination}), chúng tôi hiện thực hóa đồ thị trạng thái hội chẩn \texttt{DiagnosticStateGraph} kết hợp bộ truy xuất lai Qdrant--BM25 và bộ kiểm chứng tất định, được minh họa tại Hình~\ref{fig:multi_agent_rag_flow}.

% ==============================================================================
% Hình 2: Luồng Hội Chẩn Đa Tác Tử & Truy Xuất Tri Thức Lai (Agentic RAG Flow)
% ==============================================================================
\begin{figure}[htbp]
\centering
\resizebox{\textwidth}{!}{%
\begin{tikzpicture}[node distance=1.1cm and 1.3cm]

    % Input node
    \node[mcdssBox, fill=blue!5, text width=3.1cm] (input) {
        \textbf{Đầu vào Hội chẩn}\\
        \scriptsize Sinh hiệu $\bm{x}_{\text{tab}}$, Bệnh án $T$,\\
        Xác suất $\hat{\bm{p}}$ (Top-1, Top-2)
    };

    % Node 1: Triage & Query Formulation
    \node[mcdssAgent, right=1.2cm of input, text width=3.3cm] (triage) {
        \textbf{1. Tác tử Phân loại}\\
        \scriptsize (\textit{Triage \& Query Node})\\
        Trích xuất hội chứng \& tạo truy vấn y khoa
    };

    % Node 2: Hybrid Retriever
    \node[mcdssService, right=1.2cm of triage, text width=3.5cm] (retriever) {
        \textbf{2. Truy xuất Lai}\\
        \normalfont\scriptsize (\textit{HybridGuidelineRetriever})\\
        0.65 Từ khóa Y khoa + 0.35 Vector Ngữ nghĩa
    };

    % Node 3: Diagnostic Synthesis
    \node[mcdssAgent, below=1.2cm of retriever, text width=3.5cm] (synthesis) {
        \textbf{3. Tác tử Tổng hợp}\\
        \scriptsize (\textit{Clinical Reasoning Node})\\
        Đối chiếu chẩn đoán phân biệt \& trích dẫn CPG
    };

    % Node 4: Evidence Grounder
    \node[mcdssService, left=1.2cm of synthesis, text width=3.3cm, draw=mcdssDanger] (grounder) {
        \textbf{4. Bộ Kiểm Chứng}\\
        \normalfont\scriptsize (\textit{EvidenceGrounder} $\mathcal{O}(N)$)\\
        Đối sánh chuỗi con \& tính tọa độ $[start, end]$
    };

    % Output Verified
    \node[mcdssBox, left=1.2cm of grounder, fill=teal!8, text width=3.1cm] (output) {
        \textbf{Báo cáo Đã Xác thực}\\
        \scriptsize \texttt{hallucination\_free: true}\\
        Highlight câu bệnh án \& điều khoản GINA/GOLD
    };

    % Arrows
    \draw[mcdssArrow] (input) -- (triage);
    \draw[mcdssArrow] (triage) -- node[above, font=\tiny\sffamily] {Query + Disease Filter} (retriever);
    \draw[mcdssArrow] (retriever) -- node[right, font=\tiny\sffamily] {Top-$k$ Phác đồ} (synthesis);
    \draw[mcdssArrow] (synthesis) -- node[above, font=\tiny\sffamily] {Candidate Spans} (grounder);
    \draw[mcdssArrow] (grounder) -- node[above, font=\tiny\sffamily] {100\% Khớp gốc} (output);

    % Anti-hallucination rejection loop
    \draw[mcdssDashedArrow, draw=mcdssDanger] (grounder.south) -- ++(0,-0.55) -| node[pos=0.25, below, font=\tiny\sffamily\bfseries, text=mcdssDanger] {Loại bỏ câu trích dẫn không tồn tại trong văn bản gốc (Anti-Hallucination Gate)} (synthesis.south);

\end{tikzpicture}%
}
\caption{Đồ thị trạng thái hội chẩn đa tác tử (\textit{Multi-Agent Diagnostic StateGraph}). Mọi câu trích dẫn chứng cứ do tác tử sinh ra đều phải đi qua cổng kiểm chứng tất định \texttt{EvidenceGrounder} trước khi hiển thị cho bác sĩ.}
\label{fig:multi_agent_rag_flow}
\end{figure}



\subsection{Đồ thị Trạng thái Hội chẩn 4 Nút (\texttt{DiagnosticStateGraph})}
Luồng suy luận giải trình trong module \texttt{src/agent/diagnostic\_graph.py} được điều phối qua lớp máy trạng thái hữu hạn có hướng \texttt{DiagnosticStateGraph}, truyền đối tượng trạng thái \texttt{DiagnosticGraphState} tuần tự qua 4 nút xử lý:
\begin{enumerate}[leftmargin=1.8em, itemsep=0.2em]
    \item \textbf{Nút Phân loại \& Lập truy vấn (\texttt{triage\_and\_query\_node}):} Tiếp nhận bệnh án gốc, thực thi ẩn danh hóa thông tin định danh cá nhân qua \texttt{ClinicalDeidentifier}, xác định cặp chẩn đoán cạnh tranh sát nhất (Top-1 và Top-2, ví dụ: \texttt{ASTHMA} so với \texttt{COPD} ở ca \texttt{PT-0044}), rà quét các dấu hiệu cờ đỏ cấp cứu (\textit{Red Flags}) và khởi tạo câu truy vấn y khoa chuyên biệt.
    \item \textbf{Nút Truy xuất Tri thức Lai (\texttt{hybrid\_guideline\_retrieval\_node} --- \texttt{src/agent/retriever.py}):} Thực hiện tìm kiếm kết hợp thông qua lớp \texttt{HybridGuidelineRetriever}. Kiến trúc hỗ trợ truyền hàm mã hóa nhúng tùy chỉnh (\texttt{embedding\_fn}, chẳng hạn mô hình \texttt{sentence-transformers/all-mpnet-base-v2} theo cấu hình của nhánh Modeling); ở chế độ triển khai container CPU độc lập mặc định, module sử dụng hàm chiếu đặc trưng tất định \texttt{\_encode\_text} kết hợp 16 cụm khái niệm y sinh (\textit{biomedical concept clusters}) và 48 chiều băm $n$-gram từ/ký tự chuẩn hóa $L_2$ ($\mathbb{R}^{64}$). Khi khởi tạo, kho tri thức nạp 6 phân đoạn phác đồ chuẩn cốt lõi (\texttt{DEFAULT\_GUIDELINE\_CHUNKS} đại diện cho \texttt{GOLD\_2024}, \texttt{GINA\_2024}, \texttt{ESC\_2023}, \texttt{BYT\_2020}, \texttt{ATS\_IDSA\_2019}, \texttt{SSC\_2021}) cùng toàn bộ các tệp phân mảnh mở rộng tại \texttt{data/guidelines/*.json} vào bộ sưu tập \texttt{clinical\_guidelines} của cơ sở dữ liệu vector \textbf{Qdrant} (\texttt{QdrantClient} với độ đo \texttt{Distance.COSINE}) bằng cấu trúc \texttt{PointStruct}.
    \begin{itemize}[leftmargin=1.5em, itemsep=0.15em]
        \item \textit{Đánh giá mô hình nhúng ngữ nghĩa (Chờ bổ sung):} \texttt{[TODO-TV1 (Backlog 05 Task 1.1): Điền kết quả thử nghiệm độ tương đồng ngữ nghĩa của sentence-transformers/all-mpnet-base-v2 (vector\_dim=768, similarity\_threshold=0.72) từ configs/rag\_config.yaml]}.
        \item \textit{Quy mô kho tri thức phác đồ sau phân mảnh đầy đủ (Chờ bổ sung):} \texttt{[TODO-TV2 (Backlog 05 Task 2.1 \& Backlog 07 Task 2.1): Điền tổng số chunks phác đồ thực tế (>= 100 chunks) và phân bố theo từng nguồn GINA 2024 / GOLD 2024 / ESC 2023 / Bộ Y tế sau khi chạy src/agent/knowledge\_ingest.py]}.
    \end{itemize}
    Tại thời điểm truy vấn $q$, điểm xếp hạng của mỗi phân đoạn phác đồ $d_k$ là tổ hợp lồi giữa điểm từ khóa \textbf{Okapi BM25} chuẩn hóa ($S_{\text{bm25\_norm}}$, với tham số $k_1 = 1.5, b = 0.75$) và độ tương đồng \textbf{Cosine Vector} truy xuất từ Qdrant ($S_{\text{cosine}}$):
    \begin{equation}
        S_{\text{hybrid}}(q, d_k) = \alpha \cdot S_{\text{bm25\_norm}}(q, d_k) + (1 - \alpha) \cdot S_{\text{cosine}}(q, d_k), \quad \text{với } \alpha = 0.65
    \end{equation}
    trong đó thành phần Okapi BM25 được tính theo công thức chuẩn:
    \begin{equation}
        \text{BM25}(q, d_k) = \sum_{t \in q} \ln\!\left(1 + \frac{N_{\text{doc}} - \text{df}(t) + 0.5}{\text{df}(t) + 0.5}\right) \cdot \frac{\text{tf}(t, d_k) \cdot (k_1 + 1)}{\text{tf}(t, d_k) + k_1 \cdot \left(1 - b + b \cdot \frac{|d_k|}{\text{avgdl}}\right)}
    \end{equation}
    \item \textbf{Nút Lập luận Lâm sàng (\texttt{clinical\_reasoning\_node}):} Đối chiếu các phân đoạn phác đồ Top-$k$ thu được với hồ sơ bệnh nhân để tổng hợp báo cáo giải trình chẩn đoán phân biệt ($H_1$ so với $H_2$), đề xuất cận lâm sàng cần làm thêm và trích xuất danh sách các câu bằng chứng ứng viên $\{s_1, s_2, \dots, s_m\}$.
    \begin{itemize}[leftmargin=1.5em, itemsep=0.15em]
        \item \textit{Tích hợp Prompt Chain-of-Thought Đối kháng (Chờ bổ sung):} \texttt{[TODO-TV1 (Backlog 05 Task 1.2): Điền thông số mô hình LLM và kết quả thực nghiệm của module Prompt CoT đối kháng trong src/agent/diagnostic\_agent.py]}.
    \end{itemize}
    \item \textbf{Nút Kiểm chứng Chống ảo giác (\texttt{evidence\_grounding\_node} --- \texttt{src/agent/evidence\_grounder.py}):} Đóng vai trò cổng kiểm duyệt bắt buộc (\textit{Hard Verification Gate}) trước khi xuất kết quả ra điểm cuối \texttt{POST /api/v1/agent/explain}.
\end{enumerate}

\subsection{Thuật toán Neo Bằng chứng Ký tự $\mathcal{O}(N)$ (Character-Offset Evidence Grounding)}
Để triệt tiêu hoàn toàn khả năng tác tử tự suy diễn hoặc trích dẫn sai lệch nội dung bệnh án, module \texttt{EvidenceGrounder} thực thi thuật toán đối sánh chuỗi con tất định có độ phức tạp thời gian tuyến tính $\mathcal{O}(N)$ theo độ dài văn bản gốc $N = |T|$ (Hình~\ref{fig:evidence_grounding}).

% ==============================================================================
% Hình 3: Cơ chế Neo Bằng Chứng Ký Tự (Character-Offset Evidence Grounding)
% ==============================================================================
\begin{figure}[htbp]
\centering
\resizebox{0.96\textwidth}{!}{%
\begin{tikzpicture}[node distance=0.9cm and 1.2cm]

    % Raw Note Box
    \node[mcdssBox, text width=13.5cm, align=left, fill=mcdssBgLight] (raw_note) {
        \textbf{Văn bản bệnh án gốc ($T$, độ dài $N$ ký tự --- Ví dụ Ca PT-0044):}\\[0.2em]
        \texttt{\footnotesize "... Patient is an \colorbox{yellow!35}{\textbf{18yo female presenting with acute wheezing}} and shortness of breath after exercise. History of childhood asthma, currently using \colorbox{teal!25}{\textbf{albuterol inhaler PRN}} ..."}
    };

    % Candidate Spans
    \node[mcdssAgent, below left=1.0cm and -2.2cm of raw_note, text width=5.8cm, align=left] (cand_valid) {
        \textbf{Trích dẫn ứng viên hợp lệ $s_1$:}\\
        \texttt{\scriptsize "18yo female presenting with acute wheezing"}\\
        $\rightarrow$ Khớp chính xác tại $[start=18, end=60]$
    };

    \node[mcdssBox, below right=1.0cm and -2.2cm of raw_note, text width=5.8cm, align=left, draw=mcdssDanger, fill=red!4] (cand_invalid) {
        \textbf{Trích dẫn ảo giác $s_{\text{hallucinated}}$:}\\
        \texttt{\scriptsize "Spirometry shows FEV1/FVC ratio of 0.58"}\\
        $\rightarrow$ Không tồn tại trong $T$ ($\text{find}(s) = -1$)
    };

    % Verification Engine
    \node[mcdssService, below=2.6cm of raw_note, text width=8.5cm] (verifier) {
        \textbf{Thuật toán Đối sánh Tất định \texttt{EvidenceGrounder.verify\_and\_locate()}}\\
        \normalfont\scriptsize Độ phức tạp $\mathcal{O}(N)$ $\cdot$ Chuẩn hóa khoảng trắng $\cdot$ Bất biến chữ hoa/thường
    };

    % Final UI Output
    \node[mcdssBox, below=0.9cm of verifier, text width=11.0cm, fill=teal!6] (ui_render) {
        \textbf{Payload JSON trả về cho \texttt{EvidenceViewer.tsx}:}\\
        \texttt{\scriptsize \{"start\_char": 18, "end\_char": 60, "verified": true, "clinical\_rationale": "Khởi phát trẻ tuổi + ran rít cấp"\}}
    };

    % Connections
    \draw[mcdssArrow] (raw_note.south) -- (cand_valid.north);
    \draw[mcdssArrow, draw=mcdssDanger] (raw_note.south) -- (cand_invalid.north);
    \draw[mcdssArrow] (cand_valid.south) -- node[left, font=\tiny\sffamily\bfseries, text=mcdssSecondary] {Chấp nhận (Verified)} (verifier.north west);
    \draw[mcdssDashedArrow, draw=mcdssDanger] (cand_invalid.south) -- node[right, font=\tiny\sffamily\bfseries, text=mcdssDanger] {Bác bỏ tức thì (Rejected)} (verifier.north east);
    \draw[mcdssArrow] (verifier.south) -- (ui_render.north);

\end{tikzpicture}%
}
\caption{Cơ chế kiểm chứng vị trí ký tự $\mathcal{O}(N)$ trong \texttt{EvidenceGrounder}: Chỉ những chuỗi con tồn tại nguyên vẹn trong văn bản gốc mới được cấp tọa độ $[start, end]$ để tô sáng trên giao diện bác sĩ; mọi suy diễn ngoài văn bản đều bị chặn hoàn toàn.}
\label{fig:evidence_grounding}
\end{figure}


Với mỗi câu trích dẫn ứng viên $s_j$ do tác tử đề xuất, hàm \texttt{locate\_claim\_in\_note}$(T, s_j)$ thực hiện quy trình kiểm chứng 3 tầng:
\begin{enumerate}[leftmargin=1.8em, itemsep=0.2em]
    \item \textbf{Đối sánh chính xác tuyệt đối (\textit{Exact Substring Search}):} Tìm vị trí xuất hiện đầu tiên $p = T.\text{find}(s_j)$. Nếu $p \ge 0$, trả về cặp chỉ số ký tự $[\text{start\_char}, \text{end\_char}] = [p, p + |s_j|]$.
    \item \textbf{Đối sánh bất biến chữ hoa/thường (\textit{Case-Insensitive Fallback}):} Nếu bước 1 không khớp do khác biệt viết hoa đầu câu, thực hiện tìm kiếm trên chuỗi chữ thường $p = \text{lower}(T).\text{find}(\text{lower}(s_j))$ nhưng vẫn trích xuất chuỗi gốc $T[p : p + |s_j|]$ để bảo toàn nguyên trạng văn bản bác sĩ viết.
    \item \textbf{Đối sánh chuẩn hóa khoảng trắng (\textit{Whitespace-Normalized Regex Alignment}):} Trong trường hợp văn bản bệnh án có ký tự xuống dòng hoặc nhiều dấu cách liên tiếp, thuật toán biên dịch biểu thức chính quy thoát ký tự đặc biệt $\text{pattern} = \texttt{\textbackslash s+}.\text{join}(\text{escape}(w) \text{ for } w \in s_j.\text{split}())$ và ánh xạ ngược lại vị trí $[m.\text{start}(), m.\text{end}()]$ trên văn bản gốc $T$.
\end{enumerate}

Mọi ứng viên $s_j$ không vượt qua được quy trình kiểm chứng đều bị bác bỏ tức thì và ghi nhận vào bộ đếm kiểm toán \texttt{hallucination\_filtered\_count}. Nhờ cơ chế này, mọi đối tượng \texttt{EvidenceSpan} được gửi xuống giao diện \texttt{EvidenceViewer.tsx} đều thỏa mãn bất biến toán học:
\begin{equation}
    0 \le \text{start\_char} < \text{end\_char} \le |T| \quad \text{và} \quad T[\text{start\_char}:\text{end\_char}] \equiv \text{text}
\end{equation}
đảm bảo độ tin cậy tuyệt đối (100\% \textit{Hallucination-Free Grounding}) khi bác sĩ rà soát hồ sơ.
